% Options for packages loaded elsewhere
\PassOptionsToPackage{unicode}{hyperref}
\PassOptionsToPackage{hyphens}{url}
\PassOptionsToPackage{dvipsnames,svgnames,x11names}{xcolor}
\documentclass[
  10pt,
]{ctexart}
\usepackage{xcolor}
\usepackage[margin=2cm]{geometry}
\usepackage{amsmath,amssymb}
\setcounter{secnumdepth}{-\maxdimen} % remove section numbering
\usepackage{iftex}
\ifPDFTeX
  \usepackage[T1]{fontenc}
  \usepackage[utf8]{inputenc}
  \usepackage{textcomp} % provide euro and other symbols
\else % if luatex or xetex
  \usepackage{unicode-math} % this also loads fontspec
  \defaultfontfeatures{Scale=MatchLowercase}
  \defaultfontfeatures[\rmfamily]{Ligatures=TeX,Scale=1}
\fi
\usepackage{lmodern}
\ifPDFTeX\else
  % xetex/luatex font selection
  \setmainfont[]{DejaVu Sans}
\fi
% Use upquote if available, for straight quotes in verbatim environments
\IfFileExists{upquote.sty}{\usepackage{upquote}}{}
\IfFileExists{microtype.sty}{% use microtype if available
  \usepackage[]{microtype}
  \UseMicrotypeSet[protrusion]{basicmath} % disable protrusion for tt fonts
}{}
\makeatletter
\@ifundefined{KOMAClassName}{% if non-KOMA class
  \IfFileExists{parskip.sty}{%
    \usepackage{parskip}
  }{% else
    \setlength{\parindent}{0pt}
    \setlength{\parskip}{6pt plus 2pt minus 1pt}}
}{% if KOMA class
  \KOMAoptions{parskip=half}}
\makeatother
\usepackage{longtable,booktabs,array}
\usepackage{caption}
\captionsetup[table]{skip=6pt}
\newcounter{none} % for unnumbered tables
\usepackage{calc} % for calculating minipage widths
% Correct order of tables after \paragraph or \subparagraph
\usepackage{etoolbox}
\makeatletter
\patchcmd\longtable{\par}{\if@noskipsec\mbox{}\fi\par}{}{}
\makeatother
% Allow footnotes in longtable head/foot
\IfFileExists{footnotehyper.sty}{\usepackage{footnotehyper}}{\usepackage{footnote}}
\makesavenoteenv{longtable}
\setlength{\emergencystretch}{3em} % prevent overfull lines
\providecommand{\tightlist}{%
  \setlength{\itemsep}{0pt}\setlength{\parskip}{0pt}}
\usepackage{bookmark}
\IfFileExists{xurl.sty}{\usepackage{xurl}}{} % add URL line breaks if available
\urlstyle{same}
% fallback for those not using the hyperref driver hyperxmp:
\makeatletter
\@ifundefined{xmpquote}{\newcommand{\xmpquote}[1]{#1}}{}
\makeatother
\hypersetup{
  colorlinks=true,
  linkcolor={Maroon},
  filecolor={Maroon},
  citecolor={Blue},
  urlcolor={Blue},
  pdfcreator={LaTeX via pandoc}}

\author{}
\date{}

\usepackage{pdflscape,fvextra}
\setlength{\tabcolsep}{3pt}
\begin{document}

\section{多 Agent
文献审阅：来源与药物指向预筛}\label{ux591a-agent-ux6587ux732eux5ba1ux9605ux6765ux6e90ux4e0eux836fux7269ux6307ux5411ux9884ux7b5b}

2026-09-29 在冻结的 \texttt{multi\_\allowbreak{}agent\_\allowbreak{}review\_\allowbreak{}v1.\allowbreak{}json}
上做事后审计。审计不重跑模型，也不改写 v1 的 10 个分级。它用 PubMed
EFetch 重新获取 v1
引用的摘要，检查此前逐字引用校验没有覆盖的两个问题：勘误是否被当作研究证据，以及标为
\texttt{candidate} 的论断所引句子是否直接提及候选药的确切名称。

运行：\texttt{python\ -m\ scripts.\allowbreak{}audit\_\allowbreak{}multi\_\allowbreak{}agent\_\allowbreak{}scope}。脚本以
\texttt{TraceRecorder} 记录输入哈希、每批 PMID
和输出哈希；详细逐条结果保存在被 Git 忽略的
\texttt{artifacts/\allowbreak{}evidence\_\allowbreak{}scope\_\allowbreak{}audit/\allowbreak{}}。未来的新审阅运行会在引用校验后应用相同的预筛，将未通过项从自动分级证据中暂时移出并保留原因，供人工核对。

{\def\LTcaptype{none} % do not increment counter
\begin{longtable}[]{@{}>{\raggedright\arraybackslash}p{0.4500\textwidth}>{\raggedright\arraybackslash}p{0.4500\textwidth}@{}}
\toprule\noalign{}
结果 & 条数 \\
\midrule\noalign{}
\endhead
\bottomrule\noalign{}
\endlastfoot
v1 已通过逐字引用校验的论断 & 65 \\
需要进一步核对 & 23 \\
其中：勘误记录，不可作原始研究证据 & 5 \\
其中：候选药确切名称未出现在整篇标题/摘要 & 4 \\
其中：名称出现在记录中，但未出现在所引句子 & 14 \\
\end{longtable}
}

此前人工指出的 PMID \texttt{41376996} 勘误在 5 条引用中被发现；PMID
\texttt{33249629} 的柚皮素 A549 结果被标为 fluticasone
候选级反证，也被新规则列入待核对。后者说明``药名在摘要某处出现''依然不足以证明所引句子讲的是该药。

**边界：**23 条是待核对数，不能当作 23
条错误或``语义准确率''。药物可以用缩写、商品名、盐型或代词出现；所引句子不重复全名时也可能有效。相反，通过预筛的引文仍可能错误解读终点、研究对象或作用方向。完整的论断---引文蕴含判断仍需独立标注或人工审阅；该回顾审计不是独立测试集，也不构成疗效证据。PubMed
记录是本次重新获取的快照，未来记录变化应通过新审计的 trace 核查。

\subsection{23
条摘要级复核（事后，不改冻结版）}\label{23-ux6761ux6458ux8981ux7ea7ux590dux6838ux4e8bux540eux4e0dux6539ux51bbux7ed3ux7248}

按冻结版 claim
顺序建立了\href{../benchmark/results/evidence_scope_adjudication_v1.json}{逐条裁定账本}，并以\href{../benchmark/results/evidence_scope_adjudication_audit.json}{机器核验摘要}记录输入哈希与决策计数。审阅使用当前
PubMed 摘要快照；包含 12 篇记录、23
条原先被预筛标记的论断。完整摘要仅存于本机被 Git 忽略的
\texttt{artifacts/\allowbreak{}evidence\_\allowbreak{}scope\_\allowbreak{}audit/\allowbreak{}review\_\allowbreak{}pack\_\allowbreak{}17661de365d843\allowbreak{}16b22c11d096a5\allowbreak{}c006.\allowbreak{}json}；构建、抓取、复核均有哈希链
trace。运行
\texttt{python\ -m\ scripts.\allowbreak{}verify\_\allowbreak{}evidence\_\allowbreak{}scope\_\allowbreak{}adjudication\ -\/\allowbreak{}-review-pack\ \textless{}本机证据包\textgreater{}}
会校验冻结来源、预筛及证据包 trace、23 条完整性和裁定类别。

{\def\LTcaptype{none} % do not increment counter
\begin{longtable}[]{@{}>{\raggedright\arraybackslash}p{0.3000\textwidth}>{\raggedright\arraybackslash}p{0.3000\textwidth}>{\raggedright\arraybackslash}p{0.3000\textwidth}@{}}
\toprule\noalign{}
摘要级裁定 & 条数 & 含义 \\
\midrule\noalign{}
\endhead
\bottomrule\noalign{}
\endlastfoot
排除候选药物特异归因 & 4 & 原文仅称不明糖皮质激素，不能写成
hydrocortisone / fluorometholone 证据 \\
排除勘误作原始证据 & 5 &
\href{https://pubmed.ncbi.nlm.nih.gov/41376996/}{PMID 41376996}
是勘误通知 \\
排除论断不匹配、反证方向错误或误指别药 & 4 &
包括把\href{https://pubmed.ncbi.nlm.nih.gov/33249629/}{柚皮素 A549
抗癌实验}归给 fluticasone \\
保留但收窄至原文终点 & 9 & BDP、MF、FP/FTZ
等缩写可从同一摘要还原，但多数仅有药代、炎症、递送或体外数据 \\
仅作非癌种特异的安全提示 & 1 & 高剂量系统作用风险不能证明 NSCLC
联合治疗安全性 \\
\end{longtable}
}

因此 23 条并非全部错误：10 条可保留有限范围的语境，13
条不应继续作为原样候选级支持/反证。**冻结 v1 的
10/10``证据不足''分级原样保留，但不能把该分级当作经过此次证据过滤后重新计算的准确率或临床结论。**本轮为作者的事后摘要级审阅，未阅读全部正文或进行独立双人标注；未标记的另外
42 条也没有获得语义正确性的保证。

\end{document}
